% 主流程由左至右；離線資料與外部服務由上方提供輸入。
% 架構與流程圖共用配色；在 figure 群組外載入。
\usetikzlibrary{arrows.meta,shapes.geometric}
\definecolor{archLLM}{HTML}{DCE9F8}
\definecolor{archLLMd}{HTML}{2F6DB5}
\definecolor{archPro}{HTML}{E6F2E6}
\definecolor{archProd}{HTML}{3C8D4F}
\definecolor{archKB}{HTML}{FFF3D6}
\definecolor{archKBd}{HTML}{C58A12}
\definecolor{archSrc}{HTML}{FDE9DC}
\definecolor{archSrcd}{HTML}{D9622B}
\definecolor{archGray}{HTML}{4A4A4A}

\begin{tikzpicture}[font=\normalsize,>=Stealth,
  box/.style={draw=archProd,fill=archPro,line width=0.8pt,rounded corners=4pt,
    minimum width=3.0cm,minimum height=2.1cm,text width=2.8cm,align=center,inner sep=5pt},
  llm/.style={box,draw=archLLMd,fill=archLLM},
  kb/.style={box,draw=archKBd,fill=archKB},
  src/.style={box,draw=archSrcd,fill=archSrc,dashed,minimum height=1.3cm},
  badge/.style={circle,fill=archGray,text=white,font=\normalsize\bfseries,inner sep=2pt,minimum size=0.55cm},
  arr/.style={-{Stealth[length=2.4mm,width=2mm]},line width=1pt,archGray},
  fb/.style={arr,dashed,archSrcd}]
\fill[archPro!35,draw=archProd,dashed,rounded corners=7pt]
  (-1.75,4.9) rectangle (8.75,8.1);
\node[anchor=north west,text=archProd,font=\normalsize\bfseries] at (-1.5,8.0) {離線資料整備};
\node[box] (n1) at (0,6.3) {\textbf{資料蒐集}\\評論、影片\\與店家資訊};
\node[box] (n2) at (3.5,6.3) {\textbf{屬性與證據整理}\\實體、情境\\價格與時間};
\node[kb] (n3) at (7,6.3) {\textbf{資料儲存}\\圖譜、關聯表\\與向量索引};
\node[src] (u) at (0,3.15) {\textbf{使用者需求}\\情境、位置與時間};
\node[src] (tdx) at (10.5,3.15) {\textbf{TDX 資料}\\旅運規劃與公車};
\node[anchor=west,text=archGray,font=\normalsize\bfseries,align=left] at (12.6,3.15) {線上推薦\\餐廳為最終決策單位};
\node[llm] (n4) at (0,0) {\textbf{需求理解與規劃器}\\LLM 解析\\程式安排工具};
\node[box] (n5) at (3.5,0) {\textbf{評論語意召回}\\情境相符的\\候選餐廳};
\node[box] (n6) at (7,0) {\textbf{領域證據與粗篩}\\價格、營業\\與必要條件};
\node[box] (n7) at (10.5,0) {\textbf{公車交通分析}\\步行、候車\\乘車與轉乘};
\node[box] (n8) at (14,0) {\textbf{餐廳篩選與排序}\\三值查核\\情境與交通};
\node[llm] (n9) at (17.5,0) {\textbf{輸出與互動}\\餐廳、評論依據\\交通與未知項};
\foreach \i in {1,...,9}{\node[badge] at (n\i.north west) {\i};}
\draw[arr] (n1)--(n2);
\draw[arr] (n2)--(n3);
\draw[arr] (u)--(n4);
\foreach \a/\b in {4/5,5/6,6/7,7/8,8/9}{\draw[arr] (n\a.east)--(n\b.west);}
\draw[arr] (n3.south)--(7,4.45)--(3.5,4.45)--(n5.north);
\draw[arr] (n3.south)--(n6.north);
\draw[arr,archSrcd] (tdx.south)--(n7.north);
\draw[fb] (n9.south)--(17.5,-2.0)--(0,-2.0)--(n4.south);
\node[fill=white,inner sep=3pt,text=archSrcd,font=\small] at (8.75,-2.0) {使用者修改需求後重新查詢};
\begin{scope}[shift={(0,-3.3)}]
 \draw[archLLMd,fill=archLLM] (0,0) rectangle (0.45,0.3);
 \node[anchor=west,font=\small] at (0.6,0.15) {語言模型與互動};
 \draw[archProd,fill=archPro] (5.3,0) rectangle (5.75,0.3);
 \node[anchor=west,font=\small] at (5.9,0.15) {程式分析};
 \draw[archKBd,fill=archKB] (9.5,0) rectangle (9.95,0.3);
 \node[anchor=west,font=\small] at (10.1,0.15) {資料儲存};
\end{scope}
\end{tikzpicture}
